\documentclass[12pt]{article}
\usepackage{amssymb}
\usepackage{amsthm}
\usepackage{amsgen}
\usepackage{amsfonts}
\usepackage{amsbsy}
\usepackage[dvips]{graphicx}
%\usepackage{array, rotating} 

\newtheorem{theorem}{Theorem}[section] 
\newtheorem{corollary}[theorem]{Corollary} 
\newtheorem{conjecture}[theorem]{Conjecture} 
\newtheorem{lemma}[theorem]{Lemma} 
\newtheorem{proposition}[theorem]{Proposition} 
\newtheorem{definition}[theorem]{Definition} 
\newtheorem{example}[theorem]{Example} 
\newtheorem{axiom}{Axiom} 
\newtheorem{remark}[theorem]{Remark} 
\newtheorem{question}[theorem]{Question} 
\newtheorem{magma}[theorem]{MAGMA} 
\def\pf{{\bf proof}:\ }
\def\mod{\ {\rm mod}\ }
\def\qed{$\Box$}
\def\Spec{{\rm Spec}}
\def\fff{\mathbb{F}}
\def\aaa{\mathbb{A}}
\def\A{\mathbb{A}}
\def\ppp{\mathbb{P}}
\def\P{\mathbb{P}}
\def\qqq{\mathbb{Q}}
\def\Q{\mathbb{Q}}
\def\rrr{\mathbb{R}}
\def\R{\mathbb{R}}
\def\ccc{\mathbb{C}}
\def\cccF{F}
\def\zzz{\mathbb{Z}}
\def\Z{\mathbb{Z}}
\def\ggg{\mathbb{G}}
\def\hhh{{\cal H}}
\def\eee{{\cal E}}
\def\ra{{\rightarrow}}
\newcommand{\hra}{\hookrightarrow}


\begin{document}

\author{
David Joyner\thanks{Math Dept, USNA, wdj\@@usna.edu, supported by an NSA
Math Sciences grant. }
}
\title{Toric codes over finite fields}
\date{9-1-2003}
\maketitle

\begin{abstract}

In this note, a class of error-correcting codes is
associated to a toric variety associated to a fan defined
over a finite field $\fff_q$, analogous
to the class of Goppa codes associated to a curve.
For such a ``toric code'' satisfying certain
additional conditions, we present an efficient
list decoding algorithm for the dual code.
Many examples are given. For small $q$, many of these 
codes have parameters beating the Gilbert-Varshamov bound.
In fact, using toric codes, we construct a
$(n,k,d)=(49,11,28)$ code over $\fff_8$, which is
better than any other known code listed in 
Brouwer's tables \cite{B} for that $n$ and $k$.
We give upper and conjectural lower bounds on the minimum
distance. The upper bounds are known to be sharp in some cases.
We conclude with a discussion of some decoding methods.

\end{abstract}


\vskip .3in


\section{Introduction}

Some of the constructions discussed here are implemented in 
\cite{MAGMA} and \cite{GAP}, using the {\bf toric}
package\footnote{GUAVA version 1.7 or above contains all the
toric codes commands needed to do the constructions described here
\cite{GUAVA}.} \cite{J}.

Let $\fff=\fff_q$ denote the finite field with $q$ elements.
A {\bf $q$-ary code} is, for us, a 
a short exact sequence of vector spaces

\begin{equation}
\label{eqn:1}
0 \rightarrow 
C \stackrel{\gamma}{\rightarrow}
\fff^n \stackrel{\theta}{\rightarrow}
\fff^{n-k} \rightarrow
0.
\end{equation}
We identify $C$ with its image under $\gamma$. 
If $\fff^n$ is given the usual standard
vector space basis then the matrix of
$\gamma$ is the generating matrix of $C$
and the matrix of $\theta$ is the parity check matrix 
of $C$.
The integer $n$ is called the
{\bf length} of $C$. The dimension of $C$ is 
$k={\rm dim}_{\fff_q}(C)$. A {\bf codeword} is an
element of $C$. Since the sequence (\ref{eqn:1}) is exact, a
vector $v\in \fff^n$ is a codeword if and only if 
$\theta(v)=0$. 
The {\bf Hamming distance} between two vectors is the number of
coordinates where they differ. This metric stratifies $C$ into
subsets of codewords having the same Hamming distance from 
the origin, or ``weight''. The smallest non-zero weight
which occurs is called the {\bf minimum distance},
denoted $d=d(C)$. 

In this paper, we shall use toric varieties to construct a
code which has, for a particular $n$, $k$ and $q$, a
value of $d$ larger than that known previously.

Be begin with a related example.

\begin{example}
\label{example:RM}
Let $R=\fff [x_1,...,x_m]_{\ell}$ denote the
vector space of polynomials in $m$ variables over
$\fff$ having degree $\leq \ell$. Let
$\fff^m=\{p_1,...,p_n\}$ be some indexing of all
the points in $\fff^m$, so that $n=q^m$. Let $C$ denote the 
code
\[
C=\{(f(p_1),...,f(p_n))\ |\ f\in R\}.
\]
This code is called the {\bf Reed-Muller code},
denoted ${\cal R}(\ell,m,q)$, investigated in this
polynomial form in \cite{DGM}.

In general, the length $n$, dimension $k$, and minimum distance $d$
are known exactly for RM codes:
If $\ell=a(q-1)+b$, where $1\leq b\leq q-1$, and if
$0\leq \ell\leq m$, then
$n=q^m$, 
\[
k=
\sum_{i=0}^\ell\sum_{j=0}^{[i/q]}
(-1)^j
\left(
\begin{array}{c}
m\\ j
\end{array}
\right)
\left(
\begin{array}{c}
m-1+i-qj\\ 
m-1
\end{array}
\right),
\]
and $d=(q-b)q^{m-a-1}$.
For this, see for example \cite{DGM} or
\cite{MS}, ch 13, \S 3. It is worth noting that 
there are ``easily'' constructed functions in $R$ of minimum weight
(\cite{DGM} or page 920 in \cite{HCT}).

For example, let $m=2$, $\ell = 5$, and $q$ is ``large''. Then a basis for
$R$ is given by the monomials
$\{x_1^{e_1}x_2^{e_2}\ |\ 0\leq e_1+e_2\leq 5\}$.
We can plot the exponents $(e_1,e_2)$ of these monomials.
It is a triangle with vertices $(0,0)$, $(5,0)$,
$(0,5)$. The number of lattice points in this triangle 
is the dimension of $C$.
\end{example}

The ``toric codes'' discussed in this paper are all
subcodes of a Reed-Muller code.
They are associated to a complete non-singular $r$-dimensional 
toric variety $X$,
provided with a dense torus $T\hookrightarrow X$, defined
over $\fff_q$, and a $T$-invariant divisor $G$. 
The analogy with the above example is obtained by replacing the triangle
by a convex polytope in $\rrr^r$ and replacing 
all the points in $\fff_q^m$ by a subset of the 
$\fff_q$-rational points in $X$ (often times but not always by
those $\fff_q$-rational points of $T$, $T(\fff_q)$). 
The polytope is associated to the divisor $G$
(see \cite{E}, \cite{F}, \cite{JV}, for examples).

\section{Hansen codes}
\label{sec:hansen}

We recall briefly codes associated to a toric surface,
constructed by J. Hansen \cite{H}.
As a general matter, all varieties are non-singular, unless
stated otherwise. When there is no danger, we often identify
a variety $X$ defined over an algebraically closed field
$k$ with its set of rational points $X(k)$.

Let $\fff=\fff_q$ and let 
$\overline{\fff}$ denote a separable algebraic closure. 
Let $L$ be a lattice in $\qqq^2$ generated by $v_1, v_2\in \zzz^2$,
$P$ a polytope in $\qqq^2$,
and $X(P)$ the associated toric surface. Let $L=\zzz^2$.

There is a dense embedding of $GL(1)\times GL(1)$ into $X(P)$
given as follows.
Let $T_L=Hom_{\zzz}(L,GL(1))$ (which is $\cong GL(1)\times GL(1)$
by sending $t=(t_1,t_2)$ to $m_1v_1+m_2v_2\longmapsto
e(\ell)(t)=t_1^{m_1}t_2^{m_2}$)
and let $e(\ell):T_L\rightarrow \overline{\fff}$
be defined by $e(\ell)(t)=t(\ell)$, for $t\in T_L$ and
$\ell\in L$.

Impose an ordering on the set $T_L(\fff )$ (changing the ordering
leads to an equivalent code).
Define the code $C=C_P\subset F^n$ to be the linear code generated
by the vectors
\begin{equation}
\label{eqn:code}
B=\{(e(\ell)(t))_{t\in T_L(F)}\ |\ \ell \in L\cap P\},
\end{equation}
where $n=(q-1)^2$. In some special cases, the dimension of $C$ 
and an estimate of its minimum distance can at least be 
conjectured explicitly (see Conjecture \ref{conj:shansen} below).

J. Hansen (\cite{H}, \cite{H3})
gives\footnote{In fact, \cite{H} contains only upper bounds on $d$. 
Explicit computations using the {\bf toric} package suggested
that these upper bounds were attained. This was then proven
in \cite{H3}.} 
the minimum distance $d$ of 
such codes in the cases:

(a) $P$ is an isoceles triangle with vertices $(0,0)$,$(a,a)$,$(0,2a)$,

(b) $P$ is an isoceles triangle with vertices $(0,0)$,$(a,0)$,$(0,a)$, or

(c) $P$ is a rectangle with vertices $(0,0)$,$(a,0)$,$(0,b)$,$(a,b)$,

(d) $P$ is a trapazoid with vertices $(0,0)$, $(a,0)$, $(0,b)$, $(a,b+am)$,
where $m>0$.

\noindent
provided $q$ is ``sufficient large''. His precise result is 
recalled below.

\begin{theorem} (Hansen)
\label{thrm:hansen}
Let $a$, $b$ be positive integers. Let $P$ be the polytope 
defined in (a)-(d) above. 

\begin{itemize}
\item[(a)] Assume $q>2a+1$. The code $C=C_P$ has
\[
n=(q-1)^2,\ \ k=(a+1)^2,\ \ d= n-2a(q-1).
\]

\item[(b)] Assume $q>a+1$. The code $C=C_P$ has
\[
n=(q-1)^2,\ \ k=(a+1)(a+2)/2,\ \ d= n-a(q-1).
\]

\item[(c)] Assume $q>{\rm max}(a,b)+1$. The code $C=C_P$ has
\[
n=(q-1)^2,\ \ k=(a+1)(b+1),\ \ d= n-a(q-1)-b(q-1)+ab.
\]

\item[(d)] Assume $q>{\rm max}(a,b,b+am)+1$. The code $C=C_P$ has
\[
n=(q-1)^2,\ \ k=(a+1)(b+1)+m\frac{a(a+1)}{2},\ \ 
\]
\[
d= {\rm min}((q-a-1)(q-b-1),(q-1)(q-b-am-1)).
\]

\end{itemize}
\end{theorem}

The bounds on $q$ are best possible.

\begin{remark}
For the Hansen code of type (b), it might be worth
comparing these values of $d$ with the lower bound obtained by
puncturing a Reed-Muller code at all the points of
$\fff^2-T(\fff)$. Let $a$ be as in Hansen's theorem above.
If we consider the code ${\cal R}(a,2,q)$,
where $1<a<q/2$ and we let $d_{RM}$ denote it's minimum distance
then $d_{RM}=(q-a)q$, according to Example \ref{example:RM}
above. There are $2q-1$ points to puncture in the 
Reed-Muller code to get the Hansen code of type (b), 
so we get $d\geq d_{RM}-2q+1$. Although for $q$ large
and $a$ fixed, these are both of order $q^2$, in general
Hansen's value is much better. For example, if $q=9$ and $a=3$
then $d=40$ but $ d_{RM}-2q+1=10$.
\end{remark}

We give examples of Hansen's theorem stated above
using the {\bf toric} package \cite{J}.

\begin{example}

Part (b).

\begin{center}
\begin{tabular}{|c|c|c|c|c|c|}\hline
$q$ & $a$ & $n$ & $k$ & $d$ & using \cite{B}\\ \hline
5 & 2 &  16 & 6 & 8  &Best known\\
5 & 4 &  16 & 13 & 3  & Best possible  \\
5 & 5 &  16 & 15 & 2  & MDS\\

7 & 2 &  36 & 6 & 24  & $d=25$ is best known\\ 
7 & 3 &  36 & 10 & 18  & $d=19$ is best known \\ 

8 & 2 &  49 & 6 & 35  & $d=36$ is best known  \\ 

\hline
\end{tabular}

\end{center}

\end{example}

\section{Construction}
\label{sec:toric}

This section gives a construction which is a little more 
general than Hansen's construction recalled in
\S \ref{sec:hansen}, though it still
falls in the framework of the general class of
codes constructed in \cite{Han2}, \cite{FL}.

Again, we shall focus on the case of a toric surface,
for simplicity.

Let $M\cong \zzz^r$ be a lattice in $V=M\otimes \qqq\cong \qqq^r$
and let $N\cong \zzz^r$ denote its dual.
Let $\Delta$ be a fan (of rational cones, with
respect to $M$) in $V$ and let $X=X(\Delta)$ denote the
toric variety associated to $\Delta$. 
Let $T$ denote a dense torus in $X$.

\begin{remark}
Suppose for the moment $r=2$.
Let $\tau_1$, ..., $\tau_s$ denote the edges
of $\Delta$, ordered in a counterclockwise fashion, 
and let $v_i\in N$ denote the 
lattice point closest to the origin in $\tau_i$.
We write $v_{s+1}=v_1$.
The angle between successive vectors $v_i$, $v_{i+1}$ must be
in the range $(0,\pi)$, $1\leq i\leq s$. 
In order for $X$ to be complete, we require that
$\Delta$ contain the cones generated by each succesive pair
$v_i$, $v_{i+1}$, $1\leq i\leq s$. 
In other for $X$ to
be non-singular, we require that each successive pair
$v_i$, $v_{i+1}$, $1\leq i\leq s$, generate $N$.
(Equivalently, we require that each matrix $(v_i^t,v_{i+1}^t)$
has determinant $\pm 1$.)
\end{remark}

Let $P=P_1+...+P_n$ be a positive $1$-cycle on $X$,
where the points $P_i\in X(\fff_q)$ are distinct.
Let $G$ be a $T$-invariant divisor on $X$ which does not ``meet'' $P$,
in the sense that no element of the support of $P$ 
intersects any element in the support of $G$. We
write this as
\[
{\rm supp}(G)\cap {\rm supp}(P)=\emptyset.
\]
Some additional assumptions on $G$ and $P$ shall be made later.
Let 
\[
L(G)=\{0\}\cup \{f\in \fff_q(X)^\times \ |\ div(f)+G\geq 0\}
\]
denote the Riemann-Roch space associated to $G$.

Over $\ccc$, according to \cite{F}, \S 3.4, there is a polytope 
$P_G$ in $V$ such that $L(G)$ is spanned by the
monomials $x^a$ (in multi-index notation), for $a\in P_G\cap N$.
Over $\fff_q$, the same result holds (but remember 
that the functions in $L(G)$ are on the variety $X/\overline{\fff}$ and not on the 
set $X(\fff)$). If,
for all $a=(a_1,...,a_r)\in P_G\cap N$, we have 
$|a_i|<q$, for all $1\leq i\leq r$, then we may 
identify the polynomial functions on the variety $X/\fff$ with those on the 
set $X(\fff)$. When useful, we will assume below that $q$
is ``sufficiently large'', in this sense (depending on $G$).

Let $C=C(P,G,X,L,q)$ denote the code define by

\begin{equation}
\label{eqn:code2}
C=\{(f(P_1),...,f(P_n))\ |\ f\in L(G)\}.
\end{equation}
This is the {\bf Goppa code associated to $X$, $P$,
and $G$}. 

\begin{remark}
\label{remark:subcode}
Suppose we are given a divisor $G_1$ of $X_1=X(\Delta_1)$ and
a divisor $G_2$ of $X_2=X(\Delta_2)$, where both fans
are associated to the same lattice $M$. Suppose that
their respective polytopes satisfy $P_{G_1}\subset P_{G_2}$
and that the respective dense torii are the same,
$T_1(\fff)=T_2(\fff)$. If $S\subset T_1(\fff)$ is some 
ordered subset and $P=\sum_{p\in S}p$ then
$C(P,G_1,X_1)$ is a subspace of $C(P,G_2,X_2)$, of codimension
at most $|P_{G_2}\cap M - P_{G_1}\cap M|$.
Moreover, the minimum distance of $C(P,G_2,X_2)$
is at least as large as the minimum distance of
$C(P,G_1,X_1)$. 
\end{remark}

The dual code of (\ref{eqn:code2}) is denoted

\begin{equation}
\label{eqn:dualcode}
C^\perp = \{(c_1,...,c_n)\in \fff_q^n\ |\ \sum_{i=1}^n c_if(P_i)=0,
\forall f\in L(G)\}.
\end{equation}

It is reasonable to ask under what conditions (if any) is the map
\begin{equation}
\label{eqn:evalmap}
\begin{array}{ccc}
L(G) & \rightarrow &C\\
f & \longmapsto & (f(P_1),...,f(P_n))
\end{array}
\end{equation}
an injection? This question was investigated by S. Hansen in \cite{Han3}
and \cite{Han1} for varieties, \S 5.3, and \cite{H3} 
in the case of certain toric surfaces.


\begin{lemma} (J. Hansen)
\label{lemma:eval}
If $X=X(\Delta)$ is a non-singular toric surface and $q$ is 
sufficiently large (i.e., $q>q_0(\Delta)$, for some $q_0(\Delta)>1$)
then (\ref{eqn:evalmap}) is injective.
\end{lemma}

\pf
Though this was only proven in \S 2.3 in \cite{H3} for certain cases
(where $q_0$ was explicitly determined),
the technique applies more generally to yield the
above result.
\qed

Let $\Delta$ be a fan in a lattice $L\cong \zzz^r$. 
Let $\tau_i$ and $v_i$ be as above $1\leq i\leq s$.
Let $D_i$ denote the Weil divisor 
\begin{equation}
\label{eqn:divisor}
D_i={\rm Hom}(\tau_i^\perp\cap L^\perp,\overline{\fff}^\times),
\end{equation}
which can be regarded as the Zariski closure of an orbit of $T$. 

We now tabulate parameters $(n,k,d)$ for
several families of examples of toric codes
$C(P,G)$, where $P$ is the sum of all the
$\fff_q$-valued points in a ``dense'' torus $T$ in a toric surface $X$.
These parameters were obtained using the {\bf toric} package. 

\begin{example}
\label{example:fan7}

Let $\Delta$ be the fan generated by cones
whose edges are defined by
\[
v_1=5e_1-e_2,\ \ v_2=-e_1+5e_2,\ \ v_3=-e_1-e_2.
\]
Let $X$ be the toric variety associated to $\Delta$.
In this example, the divisor $G=5D_3$ on $X/\fff_8$ yields a 
toric code which has the parameters $(49,11,28)_8$. 
This is a new code and beats the previous best known code
(for given $n$, $k$, $q$)
having parameters $(49,11,27)_8$ \cite{B}.


\end{example}

\begin{example}
\label{example:fan1}

Let $\Delta$ be the fan generated by cones
whose edges are defined by
\[
v_1=2e_1-e_2,\ \ v_2=-e_1+2e_2,\ \ v_3=-e_1-e_2.
\]
Let $X$ be the toric variety associated to $\Delta$.

\vskip .3in

\begin{center}
\setlength{\unitlength}{.02cm}
\begin{picture}(256.00,150.00)(-120.00,0.00)
\thicklines
\put(0.00,64.00){\circle*{4}} %dot
\put(0.00,128.00){\circle*{4}} 

\put(-64.00,0.00){\circle*{4}} 
\put(0.00,-64.00){\circle*{4}} 

\put(-64.00,-64.00){\circle*{4}} 
\put(-64.00,64.00){\circle*{4}} 
\put(-64.00,128.0){\circle*{4}} 

\put(64.00,0.00){\circle*{4}} 
\put(128.00,0.00){\circle*{4}} 
\put(192.00,0.00){\circle*{4}} 
%\thinlines

\put(0.00,0.00){\vector(-1,2){64.00}} 
\put(0.00,0.00){\vector(-1,-1){64.00}} 

\put(64.00,64.00){\circle*{4}} 
\put(128.00,0.00){\circle*{4}} 
\put(0.00,128.00){\circle*{4}} 
\put(-128.00,0.00){\circle*{4}} 

\put(128.00,80.00){$\sigma_1$} 
\put(-128.00,64.00){$\sigma_2$} 
\put(50.00,-100.00){$\sigma_3$} 

\put(135.00,-75.00){$v_1$} 
\put(-75.00,135.00){$v_2$} 
\put(-75.00,-75.00){$v_3$} 

\put(0.00,-64.00){\circle*{4}} 
\put(64.00,-64.00){\circle*{4}} 
\put(128.00,-64.00){\circle*{4}} 

\put(0.00,0.00){\vector(2,-1){129.00}} %basis vector
\end{picture}
\end{center}

\vskip 1in

\noindent
Note $v_1,v_2$ generate a cone $\sigma_1$ in the 
fan $\Delta$ yet do not generate the lattice $\zzz^2$.
By the criterion on page 29 of \cite{F}, $X$ is singular.

In the notation of (\ref{eqn:divisor}), 
the divisor $G=d_1D_1+d_2D_2+d_3D_3$
is a Cartier divisor if and only if
$d_1\equiv d_2\equiv d_3\ ({\rm mod}\ 3)$ (this is
a consequence of an exercise on page 62 of \cite{F}).
A Cartier divisor is very ample if and only if
$d_1+d_2+d_3> 0$ (this follows from the
criterion in the proof of the proposition on page
68 in \cite{F}). Let 
\[
\begin{array}{c}
P_G=\{(x,y)\ |\ \langle (x,y),v_i\rangle \geq -d_i, \ \forall i\}\\
=\{(x,y)\ |\ 2x-y \geq -d_1, -x+2y \geq -d_2, -x-y \geq -d_3 
\}
\end{array}
\]
denote the polytope associated to $G$. 

Note that not all of these
divisors $G$ are Cartier. 

\begin{center}
\begin{tabular}{|c|c|c|c|c|c|c|c|}\hline
$q$ & $d_1$ & $d_2$ & $d_3$ & $n$ & $k$ & $d$ & Using \cite{B} \\ \hline
 5   &  0     &  0     &  3     &  16   &  4   &  10  & $d=11$ is best possible \\
 5   &  1     &  1     &  2     &  16   &  5   &  8  & $d=9$ is best known \\
 5   &  1     &  1     &  3     &  16   &  7   &  6  &$d=7$ is best known  \\
 5   &  1     &  1     &  4     &  16   &  10   &  4  & $d=5$ is best possible \\
 5   &  2     &  3     &  1     &  16   &  9   &  6   &Best possible   \\
 5   &  2     &  3     &  2     &  16   &  12   &  3  & $d=4$ is best possible \\
 5   &  2     &  3     &  3     &  16   &  14   &  2  & Best possible   \\
 5   &  2     &  3     &  3     &  16   &  15   &  2   & MDS\\

 7   &  1     &  2     &  0     &  36   &   3  &  30 & Best possible   \\
 7   &  1     &  2     &  2     &  36   &   7  &  21 & $d=22$ is best known  \\
 7   &  1     &  2     &  3     &  36   &   9  &  20 & Best known  \\
 7   &  3     &  0     &  0     &  36   &   4  &  27 & $d=28$ is best known  \\
8   &  0  &  0  &  3  &  49  &  4  &  40  & Best known \\
8   &  0  &  1  &  2  &  49  &  3  &  42  & Best possible \\
8   &  0  &  3  &  3  &  49  &  10  &  28  & Best known \\
8   &  0  &  3  &  4  &  49  &  12  &  26  & Best known \\

\hline
\end{tabular}

\end{center}

\end{example}


\begin{example}
\label{example:fan5}

Let for now $\Delta$ be the fan generated by 
\[
v_1=e_1,\ \ v_2=-e_1+e_2,\ \ v_3=-e_1-e_2.
\]
Let $X$ be the toric variety associated to $\Delta$.
This is singular.
The divisor $G=d_1D_1+d_2D_2+d_3D_3$ with
$d_1=0$, $d_2=0$, $d_3=a$ gives case (a) of
Hansen's code.

To create a smooth toric variety yielding Hansen's code,
let instead $\Delta$ be the refined fan generated by 
\[
v_1=e_1,\ \ v_2=-e_1+e_2,\ \ v_3=-e_1-e_2,\ \ v_4=-e_2.
\]
Let $X$ be the toric variety associated to $\Delta$.
This is now non-singular.
The divisor $G=d_1D_1+d_2D_2+d_3D_3+d_4D_4$ with
$d_1=0$, $d_2=0$, $d_3=a$, $d_4=2a$ (for example)
gives case (a) of Hansen's code.
\end{example}


\begin{example}
\label{example:fan4}

Let $\Delta$ be the fan generated by 
\[
v_1=e_1,\ \ v_2=e_2,\ \ v_3=-e_1-e_2.
\]
Let $X$ be the toric variety associated to $\Delta$.
This is non-singular.
The divisor $G=d_1D_1+d_2D_2+d_3D_3$ with
$d_1=0$, $d_2=0$, $d_3=a$ gives case (b) of
Hansen's code.
\end{example}

\begin{example}
\label{example:fan3}

Let $\Delta$ be the fan generated by 
\[
v_1=e_1,\ \ v_2=e_2,\ \ v_3=-e_1,\ \ v_4=-e_2.
\]
Let $X$ be the toric variety associated to $\Delta$.
This is non-singular.
The divisor $G=d_1D_1+d_2D_2+d_3D_3+d_4D_4$ with
$d_1=0$, $d_2=0$, $d_3=a$, $d_4=b$ gives case (c) of
Hansen's code.

\end{example}


\begin{example}
\label{example:fan2}

Let $\Delta$ be the fan generated by 
\[
v_1=e_1,\ \ v_2=-e_1+me_2\ \ (m>1\ {\rm fixed}),\ \ v_3=-e_2.
\]
Let $X$ be the toric variety associated to $\Delta$.
This surface is, except for small $m$, singular.

In the notation above, the divisor $G=d_1D_1+d_2D_2+d_3D_3$
is a Cartier divisor if and only if
$d_1 + d_2\equiv 0\ ({\rm mod}\ m)$ (this is
an exercise on page 64 of \cite{F}).

Let $\Delta'$ be the refined fan generated by 
\[
v_1=e_1,\ \ v_2=e_2,\ \ v_3=-e_1+me_2\ \ (m>1\ {\rm fixed}),
\ \ v_4=-e_2.
\]
Let $X'$ be the toric variety associated to $\Delta'$.
This surface is non-singular.
The divisor $G'=d'_1D'_1+d'_2D'_2+d'_3D'_3+d'_4D_4$
(using $D'_i$ to denote the divisor in
(\ref{eqn:divisor}) defined on $X'$)
is very ample if and only if
$d'_2+d'_4\geq 0$, $d'_1+d'_3\geq md'_1$ (this is an exercise 
on page 70 of \cite{F}).
Taking 
\[
d_1'=d_2'=0, \ \ d_3'=a,\ \ d_4'=b,
\]
and swapping the roles of $x$ and $y$, gives the Hansen code
of type (d) constructed above.

Let 
\[
\begin{array}{c}
P_G=\{(x,y)\ |\ \langle (x,y),v_i\rangle \geq -d_i, \ \forall i\}\\
=\{(x,y)\ |\ x \geq -d_1, -x+my \geq -d_2, -y \geq -d_3 
\}
\end{array}
\]
denote the polytope associated to the Weil
divisor $G=d_1D_1+d_2D_2+d_3D_3$.

The {\bf toric} package gives the following results:

$m=3$

\begin{center}
\begin{tabular}{|c|c|c|c|c|c|c|c|}\hline
$q$ & $d_1$ & $d_2$ & $d_3$ & $n$ & $k$ & $d$ & Using \cite{B} \\ \hline
5    & 0  &  2 & 2  & 16  & 11  & 3 & $d=4$ is best known  \\
5    & 0  &  2 & 3  & 16  & 15  & 2 & MDS  \\
5    & 0  &  4 & 2  & 16  & 14  & 2 & Best possible  \\
5    & 1  &  0 & 0  & 16  & 2  & 12 & $d=13$ is best possible  \\
7    & 0  &  1 & 0  & 36  & 2  & 30 & $d=31$ is best possible  \\
8    & 1  &  0 & 0  & 49  & 2  & 42 & $d=43$ is best possible  \\
9    &  0 &  1 & 0  & 64  & 2  & 56 & $d=57$ is best possible  \\
\hline
\end{tabular}
\end{center}

$m=5$ 

\begin{center}
\begin{tabular}{|c|c|c|c|c|c|c|c|}\hline
$q$ & $d_1$ & $d_2$ & $d_3$ & $n$ & $k$ & $d$ & Using \cite{B} \\ \hline
5    & 0  &  0 & 3  & 16  & 13 & 2 & $d=3$ is best possible  \\
5    & 3  &  3 & 2  & 16  & 14 & 2 & Best possible  \\
\hline
\end{tabular}
\end{center}

Indeed, the parameters arising in these cases are very similar
to each other.
\end{example}


\begin{example}
\label{example:fan6}

Let $\Delta$ be the fan generated by 
\[
v_1=2e_1-e_2,\ \ v_2=-e_1+e_2,\ \ v_3=-e_1.
\]
Let $X$ be the toric variety associated to $\Delta$.
This is non-singular.
All the divisors $G=d_1D_1+d_2D_2+d_3D_3$
are Cartier.
The divisor $G=d_1D_1+d_2D_2+d_3D_3$
is very ample if and only if $d_1+d_2+d_3>0$.

Let 
\[
\begin{array}{c}
P_G=\{(x,y)\ |\ \langle (x,y),v_i\rangle \geq -d_i, \ \forall i\}\\
=\{(x,y)\ |\ 2x-y \geq -d_1, -x+y \geq -d_2, -x \geq -d_3 
\}
\end{array}
\]
denote the polytope associated to $G$. 
If each $d_i>0$ and if $d_3>d_1+d_2$ then the volume of this polytope 
is given by
\[
{\rm vol}(P_G)=\frac{1}{2}(2d_1+d_3)(d_3-d_1-d_2).
\]

\begin{center}
\begin{tabular}{|c|c|c|c|c|c|c|c|}\hline
$q$ & $d_1$ & $d_2$ & $d_3$ & $n$ & $k$ & $d$ & Using \cite{B} \\ \hline
5    &  0     &  0     &  1     &  16   &  3   &  12 & Best possible   \\
5    &  0     &  0     &  2     &  16   &  6   &  8  & $d=9$ is best known \\
5    &  0     &  0     &  3     &  16   &  10   &  4  & $d=5$ is best known \\
5    &  0     &  0     &  4     &  16   &  13   &  3 & Best possible   \\
7    &  0     &  0     &  1     &  36   &  3   &  30  & Best possible   \\
7    &  0     &  0     &  2     &  36   &  6   &  24  & Best known \\
7    &  0     &  0     &  3     &  36   &  10   &  18  & Best known \\
7    &  0     &  0     &  4     &  36   &  15   &  12  & $d=14$ is best known \\
7    &  4     &  1     &  1     &  36   &  26   &  5  & $d=6$ is best known \\
7    &  4     &  1     &  2     &  36   &  30   &  4 & $d=5$ is best possible    \\
7    &  4     &  1     &  3     &  36   &  33   &  3 & Best possible   \\
7    &  4     &  1     &  4     &  36   &  35   &  2  & MDS \\
8    &  0     &  0     &  1     &  49   &  3   &   42  & Best possible \\
8    &  0     &  0     &  2     &  49   &  6   &   35  & $d=36$ is best known   \\
8    &  0     &  0     &  3     &  49   &  10   &  28  & Best known \\
8    &  0     &  0     &  4     &  49   &  15   &  21  & $d=23$ is best known \\
8    &  0     &  4     &  4     &  49   &  39   &  5 & $d=6$ is best known   \\
8    &  2     &  4     &  4     &  49   &  46   &  3 & best possible   \\
8    &  3     &  4     &  4     &  49   &  48   &  2 & MDS  \\
8    &  4     &  1     &  4     &  49   &  43   &  4 & $d=5$ is best possible   \\
9    &  0     &  0     &  1     &  64   &   3  &   56  & Best possible  \\
9    &  0     &  0     &  2     &  64   &   6  &   48  & $d=49$ is best known  \\
9    &  0     &  0     &  3     &  64   &   10  &  40  & $d=41$ is best known \\

\hline
\end{tabular}

\end{center}

\end{example}

\section{Estimates on the parameters}
\label{sec:estimate}

For toric codes, estimates for $d$ boils down to 
knowing how to estimate an answer to the following natural question:
given any polynomial $f\in \fff [x_1,...,x_m]$, where $\fff$ is
a finite field with $q$ elements, of degree $\ell$, how many zeroes
can $f$ have? The answer can depend on $m$, $q$, and $\ell$
but not on $f$.

We shall now estimate the parameters $n,k,d$ for 
$C=C(P,G,X)$, conjecturally. Assume $X$ is a complete toric
variety of dimension $r$. The form of these conjectures was
inspired by research of S. Hansen \cite{Han3}.
For many such varieties,
the conjectured estimates below were verified by computer 
using the {\bf toric} package.

Let $G$ be a $T$-invariant Cartier divisor on $X=X(\Delta)$,
let $P_G$ be the polytope associated to $G$ and let $M$ be 
the lattice as in \cite{F}.


\begin{conjecture}
\label{conj:shansen}
Let $C=C_L(P,G,X)$ be as in (\ref{eqn:code2}). Assume 
\begin{itemize}
\item
$X$ is a non-singular,
projective toric variety of dimension $r$, 
\item
$n$ is so large that there is an integer $N>1$ such that
\[
2N\cdot {\rm vol}(P_G)\leq n\leq 2N^2\cdot {\rm vol}(P_G).
\]
\end{itemize}
If $q$ is ``large'' then any 
$f\in L(G)=H^0(X,{\cal O}(G))$ has no more than $n$ zeros
in $X(\fff_q)$.
Consequently, 

\[
d\geq n- 2N\cdot {\rm vol}(P_G).
\]

\end{conjecture}

Here ``large'' make depend on $X$, $P$ and $G$ but not
on $f$.

Let $G$ be a $T$-invariant Cartier divisor on $X=X(\Delta)$
and let
\[
\psi_G:|\Delta|\rightarrow \rrr
\]
be the linear function associated to $G$ as in \cite{F},
\S 3.4 (this is called a {\bf virtual support
function} in \cite{E}, Definition VII.4.1).
If $\psi_G$ is strictly convex and $X$ is non-singular
then $G$ is very ample
and ${\cal O}(G)$ is generated by its global sections
(\cite{F}, \S 3.4, p. 70)\footnote{
Demazure has shown that on a complete, non-singular toric variety
a $T$-invariant divisor is ample if and only if it is very ample
(see \cite{F}, page 71).
}.
Conversely, if ${\cal O}(G)$ is generated by its global sections then
$\psi_G$ is convex and 
\[
\psi_G(v)={\rm min}_{u\in P_G\cap M}\ \langle u,v\rangle,
\]
where $P_G$ and $M$ are as above.

\begin{conjecture}
\label{conj:main}
Let $C_L=C_L(G,P,X)$ be as in (\ref{eqn:code2}). Assume 
\begin{itemize}
\item
$X$ is a non-singular,
projective toric variety of dimension $r$, 
\item
$\psi_G$ is strictly convex,
\item
and  
\[
{\rm deg}(P)>{\rm deg}(G^r).
\]
\end{itemize}
If $q$ is ``large'' then any 
$f\in L(G)=H^0(X,{\cal O}(G))$ has no more than $n={\rm deg}(P)$ zeros
in $X(\fff_q)$.
Consequently, 

\[
k\geq {\rm dim}\, H^0(X,{\cal O}(G))=|P_G\cap M|,
\]
and
\[
d\geq n-r!\cdot |P_G\cap M|.
\]
Moreover, if $n>r! |P_G\cap M|$ then 
${\rm dim}\, H^0(X,{\cal O}(G))=|P_G\cap M|$.
\end{conjecture}

\begin{remark}
This can be false if $X$ is singular (as examples show).
\end{remark}

\begin{definition}
Let $T\subset X$ be the dense torus of a toric variety $X$
as in the above conjecture.
Suppose that $P=\sum_{p\in T(\fff)}p$ and $G$ is a sum of Weil divisors.
In this case, we call $C(P,G,X)$ a {\bf standard toric code}.
\end{definition}

Suppose that the polygon $P_G$
is contained in the simplex 
$S_\ell=\{(x_1,...,x_r)\ |\ x_i\geq 0,
x+1+...+x_r\leq \ell\}$, for some integer $\ell$. 
Suppose that the polygon $P_G$
contains the simplex $S_k=\{(x_1,...,x_r)\ |\ x_i\geq 0,
x+1+...+x_r\leq k\}$, for some integer $k$. 
In the notation of Remark \ref{remark:subcode},
let $\Delta$ denote the fan associated to $X$,
$\Delta_1$ denote the fan associated to the simplex $S_k$,
let $G_1$ be a divisor in $X_1=X(\Delta_1)$ for which
$P_{G_1}=S_k$, let
$\Delta_2$ denote the fan associated to the simplex $S_\ell$,
and let $G_2$ be a divisor in $X_2=X(\Delta_2)$ for which
$P_{G_2}=S_\ell$. The codes $C(P,G_i,X_i)$ 
generalize the Hansen codes of type (b) in \S 2.

\begin{theorem}
\label{thrm:subcode2}
Under the above conditions for a standard toric code,
we have

\begin{enumerate}
\item
$C(P,G_1,X_1)$ is a subspace of $C(P,G,X)$, of codimension
at most $|P_{G_1} - S_k|$,

\item
$C(P,G,X)$ is a subspace of $C(P,G_2,X_2)$, of codimension
at most $|S_\ell - P_G|$,

\item
$C(P,G_1,X_1)$ is the subcode obtained
from ${\cal R}(k,r,q)$ (in Example \ref{example:RM})
by puncturing at the points of $\fff^r -T(\fff) $.

\item
$C(P,G_2,X_2)$ is the subcode obtained
from ${\cal R}(\ell,r,q)$ (in Example \ref{example:RM})
by puncturing at the points of $\fff^r -T(\fff) $.

\item
The minimum distance of the toric code
$C(P,G_i,X_i)$ is at least $(q-a)q^{r-1}-r(q-1)^{r-1}-1$,
where $a=k$ if $i=1$ and $a=\ell$ if $i=2$.

\item
The minimum distance of the toric code
$C(P,G,X)$ is at least $(q-\ell)q^{r-1}-r(q-1)^{r-1}-1$
and at most $(q-k)q^{r-1}$.
\end{enumerate}
\end{theorem}

\begin{remark}
In fact, we only need that {\it some translate} of
$P_G$ is contained in the simplex 
$S_\ell$ and some translate of $P_G$ contains $S_k$.
\end{remark}

\pf
Must of this follows from Remark \ref{remark:subcode}.
Each puncture decreases the minimum distance at most by 
$1$, so from the equation for the minimum distance
of the Reed-Muller code $d_{RM}$, we obtain a lower
bound for the minimum distance of such the toric code
$C(P,G_i,X_i)$, for $i=1,2$. Since raising the dimension of the 
code (while preserving the length), decreases the 
minimum distance, the minimum distance of 
$C(P,G,X)$ is at least that of $C(P,G_1,X_1)$ and
at most that of $C(P,G_2,X_2)$.
\qed

\begin{example}
Fix $r$ satisfying $1<r<q-1$. Let $C$ be a Hansen code of type (c) with parameters
$(a,b)=(r,r)$, let $C_1$ be a Hansen code of type (b) with parameter
$a=r$, and  let $C_2$ be a Hansen code of type (b) with parameter
$a=2r$. We have 

\[
(q-1)^2-2r(q-1) = min.dist. (C_2)
 \leq min.dist. (C)=(q-1)^2-2r(q-1)+r^2\leq min.dist. (C_1)=(q-1)^2-r(q-1).
\]
\end{example}

We conclude with an upper bound on the minumum distance.
Let $M^+=\{m\in M\ |\ m_i\geq 0\}$.

\begin{theorem}
Let $C=C_L(G,P,X)$ be as in (\ref{eqn:code2})
and assume $q$ is ``sufficiently large'' as in Lemma
\ref{lemma:eval}. Let $H\subset P_G\cap M^+$
be a line 
\newline
$\{(e_1,...,e_{j-1},x,e_{j+1},...,e_r)\ |\ 0\leq x\leq h\}$,
for some integer $h>0$.
Then $d\leq n-h$.

\end{theorem}

\begin{remark}
In particular, let $r=2$ and let $h>0$ denote the 
largest integer such that $\{(0,0),...,(0,h)\}\subset P_G$.
Then $d\leq n-h$.
This generalizes some results in Hansen's theorem stated above
(which are known to be sharp).

\end{remark}

\pf
Let $h={\rm max}_i(d_i)$. Let $f\in L(G)$ denote the function
\[
(x_j-a_1)...(x_j-a_h)\prod_{i\not= j}x_i^{e_i}\in L(G),
\]
where $a_i\in \fff_q^\times$ are distinct.
This vanishes at $h$ points, so $d\leq n-h$.
\qed


\section{Decoding algorithm}
\label{sec:decode}

A number of decoding algorithms are possible for toric codes.

\begin{itemize}
\item
Use a more general algorithm for
decoding multi-dimensional cyclic
codes. We refer to J. Little's papers \cite{L1}, \cite{L2} 
and the references cited there. Most or all of these
require Gr\"obner bases techniques, which become impractical
(for example) if $q$ is large or the number of errors is large.

\item
Use a more general algorithm for
decoding affine variety codes \cite{FL}. This too
requires Gr\"obner bases techniques.

\item
Since these are subcodes of a Reed-Muller code,
use a majority logic decoding algorithm valid
for the larger Reed-Muller code. 
\end{itemize}
The first two of these can be impractical\footnote{However, much of the
Gr\"obner basis computation is ``one-time work'', so it might not be 
as impractical as it seems, depending on the application.}
for toric codes since 
it is can be necessary to take $q$ to be relatively large to get
good error correction capabilities. The third does not always
apply (see ch. 13, \S 7 of \cite{MS}. 

In this section, we discuss an new algorithm.
The algorithm here will decode $C=C(P,G,X)^\perp$ when
the divisor $G$ and the number of errors satisfy certain
conditions (see (A)-(C) below).
What it lacks in optimality, it makes up for in simplicity.
It also avoids Gr\"obner bases and basically works equally fast 
for large or small $q$.

Let $X=X(\Delta)$ be a toric variety associated to a fan $\Delta$
with edges $\tau_1$, ..., $\tau_s$. Let $D_i$ denote the $T$-Weil
divisors as above.

Let $C$ be as in (\ref{eqn:dualcode}).
Let $r=c+e$ be a received vector, $c\in C$ a codeword
and $e$ the error vector of smallest weight. Thus all
$f\in L(G)$ satisfy $\sum_{i=1}^n c_if(P_i)=0$. Let
\[
I=\{i\ |\ e_i\not= 0\}.
\]

We shall follow \cite{S}, pages 217-218, to (a) create an
error locator function for $r$ and $C$, and (b) solve for 
$e_i$, $i\in I$. The method generally only gives a list
of several vectors, one of which is the true error
vector $e$ - it only yields one vector (namely, $e$)
under a certain technical condition (see
Proposition \ref{prop:errors} for details).

Let $G'$ be another divisor of $X$. We shall make some
assumptions on $G$ and $G'$ as we go along. Let
\[
L(G')={\rm span}\{f_1,...,f_\ell\},
\]
\[
L(G-G')={\rm span}\{g_1,...,g_k\},
\]
\[
L(G)={\rm span}\{h_1,...,h_m\}.
\]

We assume 
\begin{itemize}
\item[(A)]
that these spaces $L(G),L(G'),L(G-G')$ are non-zero,
\item[(B)]
that there are $T$-invariant divisors $D_i$ for which 
(a) each $P_i$ is contained in a $D_i$, $i\in I$,
(b) $P_i\notin D_j$, for $i\not= j$ and $i,j\in I$,
(c) $L(G'-\sum_{i\in I}D_i)$ is non-zero,
(d) for any fixed $i_0\in I$,
$L(G-G'-\sum_{i\in I}D_i)$ is non-empty and {\it properly}
contained in 
$L(G-G'-\sum_{i\in I,i\not= i_0}D_i)$.
\item[(C)]
$d> {\rm deg}(G')$.

\end{itemize}
Assumption (B) above is very strong and forces us
to assume that $|I|\leq s$. Roughly speaking, $G$ is a ``large'' 
divisor and the number of errors is ``small''.
(Equivalently, $G$ is ``large'' and $\Delta$ has a large number of
edges.) Assumption (C) can be removed (or at least replaced by a
much weaker condition) in the case when one of the conjectures
of the previous section is known or when the hypothesis to
Theorem \ref{thrm:subcode2} applies. 

Let
\[
[r,\phi ]=\sum_{i=1}^n r_i\phi (P_i),
\]
and let $x_1=a_1$, ..., $x_\ell=a_\ell$ be a non-trivial solution 
to the system of $k$ simultaneous equations
\begin{equation}
\label{eqn:coeffs1}
\sum_{j=1}^\ell [r,f_jg_i]x_j=0,
\ \ \ \ \ \ \ 1\leq i\leq k.
\end{equation} 
(We show such a solution exists below.)
Let $f=a_1f_1+...+a_\ell f_\ell$.

\begin{proposition}
\label{prop:errorlocator}
Under the assumptions above, 
$f$ is an error locator function. In other words,
$f(P_i)=0$ for all $i\in I$.
\end{proposition}

\pf
We have $L(G'-\sum_{i\in I}D_i)$ is non-zero, by assumption.
Choose a non-zero $z\in L(G'-\sum_{i\in I}D_i)\subset L(G')$
and write $z=\gamma_1f_1+..+\gamma_\ell f_\ell$, for some
$\gamma_i\in \fff_q$. Then $zg_j\in L(G)$, for all
$1\leq j\leq k$. If $r=(r_1,...,r_\ell)$ then 
\[
[r,zg_j]=\sum_{i=1}^\ell [r,f_ig_j]\gamma_i.
\]
On the other hand, $[r,zg_j]=[c+e,zg_j]=[e,zg_j]$, so
\[
[r,zg_j]=[e,zg_j]
= \sum_{v=1}^n e_v f_i(P_v)g_j(P_v).
\]
But $e_v=0$ for $v\notin I$ and $z(P_i)=0$ for
$v\in I$, so $[r,zg_j]=0$. This implies that
$(\gamma_1,...,\gamma_\ell)$ is a non-trivial solution to 
(\ref{eqn:coeffs1}).

Now suppose that there is an $i_0\in I$ for which 
$f(P_{i_0})\not= 0$. By assumption (B,d),
there is an $h\in L(G-G')$ such that
$h(P_{i_0})\not= 0$ and $h(P_i)=0$ for all
other $i\in I$. Then
\[
[r,fh]=\sum_{v=1}^n e_v f(P_v)h(P_v)
=e_{i_0}f(P_{i_0})h(P_{i_0})\not= 0.
\]
On the other hand, $h$ is a linear combination 
of the $g_j$'s. So this contradicts the 
equation $[r,zg_j]=0$ we obtained earlier.
\qed

Next, we must determine the $e_i$, $i\in I$.
Let 
\[
N(f)=\{i \ |\ 1\leq i\leq n, f(P_i)=0\}.
\]
By the proposition, $I\subset N(f)$. 

\begin{proposition}
\label{prop:errors}
The $b_i=e_i$, for $i\in N(f)$, solve 
\[
\sum_{i\in N(f)} b_i h_j(P_i)=[r,h_j],\ \ \ \ \ 1\leq j\leq m.
\]
If $|N(f)|\leq {\rm deg}(G')$ then this solution is unique.
\end{proposition}

\begin{remark}
Naturally, it would be interesting to know if
the hypothesis $|N(f)|\leq {\rm deg}(G')$ can be removed.

A bound on $|N(f)|$ is known for certain $\Delta$ - 
see Theorem \ref{thrm:subcode2} and
\S 2.3 in \cite{H3}. A conjectural bound is given in the
previous section.
\end{remark}

\pf 
Since $h_j\in L(G)$ and (by the previous proposition)
$e_v=0$ for all $v\notin N(f)$, we have
\[
[r,h_j]=[e,h_j]=\sum_{v=1}^n e_v h_j(P_v)
=\sum_{v\in N(f)} e_v h_j(P_v).
\]
Thus the equation displayed in Proposition
\ref{prop:errors} has the $e_i$'s, $i\in N(f)$, as a solution.

Now we show that this solution is unique. Suppose
$b_i$, $i\in N(f)$, is another solution. Define
$b_j=0$ for $j\in \{1,2,...,n\}-N(f)$. Then
for $1\leq j\leq m$,
\[
[b,h_j]=\sum_{v=1}^n b_v h_j(P_v)
=\sum_{v\in N(f)} b_v h_j(P_v)=[r,h_j]=[e,h_j].
\]
As the $h_j$'s form a basis of $L(G)$, this implies
$b-e$ is a codeword. The weight of this codeword
satisfies 
\[
wt(b-e)\leq |N(f)|\leq deg(G')< d,
\]
by our hypothesis. This forces $b_i=e_i$,
so the solution is unique. 
\qed

The above Proposition tells us how to find the error vector, given 
$f$, $L(G)$, and the $P_i$'s, as desired.

We illustrate this algorithm with an example.

\begin{example}

Here is an example where we decode 
a received word in a toric code with $\leq 3$ errors.

Let $\Delta$ be the fan as in Example \ref{example:fan1}.
Let $X$ be the toric variety associated to $\Delta$.

Let 
\[
\begin{array}{c}
P_D=\{(x,y)\ |\ \langle (x,y),v_i\rangle \geq -d_i, \ \forall i\}\\
=\{(x,y)\ |\ 2x-y \geq -d_1, -x+2y \geq -d_2, -x-y \geq -d_3 
\}
\end{array}
\]
denote the polytope associated to the Weil
divisor $D=d_1D_1+d_2D_2+d_3D_3$,
where $D_i$ is as above. 

Let 
\[
G=10D_3,\ \ \ 
D=D_1+D_2+D_3,\ \ \ G'=2D.
\]
Then $P_G$ is a triangle in the plane 
with vertices at $(0,0)$, $(10/3,20/3)$, and $(20/3,10/3)$. 
It's area is $\frac{1}{2}bh=50/3$.

Let $T\subset X$ denote the dense torus of $X$.
In this example, the patch $U_{\sigma_1}$ is an affine variety with 
coordinates $x_1,x_2,x_3$ given by $x_1^3-x_2x_3=0$. 
The torus embedding
$T\hookrightarrow U_{\sigma_1}$ is given by
sending $(t_1,t_2)$ to 
$(x_1,x_2,x_3)=(t_1t_2,t_1t_2^2,t_1^2t_2)$.
The patch $U_{\sigma_2}$ is an affine variety with 
coordinates $y_1,y_2,y_3$ given by $y_2^2-y_1y_3=0$. 
The torus embedding
$T\hookrightarrow U_{\sigma_2}$ is given by
sending $(t_1,t_2)$ to 
$(y_1,y_2,y_3)=(t_1^{-2}t_2^{-1},t_1^{-1},t_1^{-1}t_2)$.
The patch $U_{\sigma_3}$ is an affine variety with 
coordinates $z_1,z_2,z_3$ given by $z_2^2-z_1z_3=0$. 
The torus embedding
$T\hookrightarrow U_{\sigma_3}$ is given by
sending $(t_1,t_2)$ to 
$(x_1,x_2,x_3)=(t_1^{-1}t_2^{-2},t_2^{-1},t_1t_2^{-1})$.

In the local coordinates of $U_{\sigma_1}$, the
space $L(G)$ has as a basis,
\[
\begin{array}{c}
\{h_i\}=
\{ 1, x_1 x_2, x_1^2 x_2, x_1 x_2^2, x_1^2 x_2^2, 
x_1^3 x_2^2, x_1^4 x_2^2, x_1^2 x_2^3, x_1^3 x_2^3,\\ 
x_1^4 x_2^3, x_1^5 x_2^3, x_1^6 x_2^3, x_1^2 x_2^4, 
x_1^3 x_2^4, x_1^4 x_2^4, x_1^5 x_2^4, x_1^6 x_2^4, \\
x_1^3 x_2^5, x_1^4 x_2^5, x_1^5 x_2^5, x_1^3 x_2^6, 
x_1^4 x_2^6\}\, .
\end{array}
\]
In particular, it is 22-dimensional.

The polytope $P_{G'}$ has vertices at
$(-2,-2)$, $(2,0)$, $(0,2)$.
In the local coordinates of $U_{\sigma_1}$, the
space $L(G')$ has as a basis,
\[
\{f_i\}=
\{ x_1^{-2}x_2^{-2}, x_1^{-1}, x_2^{-1}, x_1^{-1}x_2^{-1},  1, x_1, x_1^2, 
x_1x_2, x_2, x_2^2\}.
\]
In particular, it is 10-dimensional.

The polytope $P_{G-G'}$ has vertices at
$(2,2)$, $(10/3,14/3)$, $(14/3,10/3)$.
In the local coordinates of $U_{\sigma_1}$, the
space $L(G-G')$ has as a basis,
\[
\{g_i\}=
\{x_1^2 x_2^2, x_1^3 x_2^3, x_1^4 x_2^3, x_1^3 x_2^4, x_1^4 x_2^4\}.
\]
$L(G-G')$ is 5-dimensional.

Choose distinct points $P_1,...,P_n\in X(\fff_q)$ and let
\[
C=\{(c_1,...,c_n)\in \fff_q^n\ |\ \sum_{i=1}^n c_i f(P_i)=0,
\forall f\in L(G)\}.
\]

Let $P_i\in D_i$ and $P_i\notin D_j$, with
$1\leq i\not= j\leq 3$. Assume $P_4,...,P_n\notin D_i$, $i=1,2,3$.
(For example, let $\{P_4,...,P_n\}= T(\fff_q)$.)
We must choose $n>50$, assuming 
Conjecture \ref{conj:main} below, to insure condition
(C) above is satisfied\footnote{For 
example, we know
\[
\begin{array}{c}
|X(\fff_2)|=7,\ \ |X(\fff_3)|=13,\ \ |X(\fff_4)|=21,\ \\  
|X(\fff_5)|=31,\ \ |X(\fff_7)|=57,\ \  |X(\fff_8)|=73,...
\end{array}
\]
So, we may take $q\geq 7$, at least conjecturally. 
Using Theorem \ref{thrm:subcode2} and the fact that $C$ differs from a 
standard toric code $C(P.G,X)$ by $3$ punctures, we may take
$q$ so large that $(q-10)q-2q+1>6$, so $q\geq 13$ will work
unconditionally.
}.

Since $L(G-3D)$ is 1-dimensional, hence non-zero,
the hypotheses and assumptions above are all
satisfied. 

Suppose $r=(1,1,1,0,...,0)\in\fff_q^n$ is the
received vector (it has $3$ errors in the first $3$ positions).
In this case,
\[
[r,f]=f(P_1)+f(P_2)+f(P_3).
\]
Let $a_1$, $a_2$, ..., $a_{10}$ be a solution to
\[
\begin{array}{c}
a_1[r,f_1g_1]+a_2[r,f_2g_1]+...+a_{10}[r,f_{10}g_1]=0,\\
a_1[r,f_1g_2]+a_2[r,f_2g_2]+...+a_{10}[r,f_{10}g_2]=0,\\
\vdots\\
a_1[r,f_1g_{5}]+a_2[r,f_2g_{5}]+..+a_{10}[r,f_{10}g_{5}]=0.
\end{array}
\]
This determines
\[
f=\sum_{i=1}^{10} a_jf_j,
\]
and hence $N(f)$.
The equations
\[
\sum_{i\in N(f)} b_i h_j(P_i)=[r,h_j],\ \ \ \ \ 1\leq j\leq 22,
\]
clearly have a solution ($b_1=b_2=b_3=1$, the other $b_i=0$).
If $q\geq 13$ then it is unique.
\end{example}


{\bf Acknowledgements}:
I thank Soren Hansen, Johan Hansen, and Amin Shokrollahi for 
answering many questions and helpful
correspondence regarding their research.

\begin{thebibliography}{99}  


\bibitem[B]{B} W. Brouwer's tables on
Bounds on the minimum distance of linear codes,
\newline
\verb+http://www.win.tue.nl/~aeb/voorlincod.html+
\newline
See also W. Brouwer's (less up-to-date) article in 
{\bf Handbook of coding theory}, (ed. Pless et al),
North-Holland, Elsevier, (1998).

\bibitem[DGM]{DGM} P. Delsarte, J. Goethals, F. MacWilliams,
``On generalized Reed-Muller codes and their relatives,''
Information and Control \underbar{16}(1970)403-442.


\bibitem[E]{E} G. Ewald, {\bf Combinatorial convexity and algebraic geometry}, 
Springer, 1996.

\bibitem[FL]{FL} J. Fitzgerald, R. Lax, ``Decoding affine varieties
using Gr\"obner bases,'' Designs, Codes, and Cryptography
\underline{13}(1998)147-158.

\bibitem[F]{F}
W. Fulton, {\bf Introduction to toric varieties}, PUP, 1993.

\bibitem[GAP]{GAP}
The GAP~Group, {\bf GAP -- Groups, Algorithms, and Programming, 
Version 4.3}; 2000, \verb+(http://www.gap-system.org)+.

\bibitem[GUAVA]{GUAVA}
The GAP package {\bf GUAVA, Version 1.6}, 2002, 
\verb+(http://www.gap-system.org/Info4/deposit.html)+

\bibitem[HCT]{HCT} V. Pless and W. Huffman (eds.),
{\bf Handbook of coding theory}, North-Holland, 1998.

\bibitem[H1]{H} J. Hansen, ``Toric surfaces and error-correcting
codes,'' {\bf Coding theory, cryptography, and related areas},
(ed., Bachmann et al), Springer-Verlag, 1999.

\bibitem[H2]{H2} -----, ``Hirzebruch surfaces and error-correcting
codes,'' preprint, 1999. Available at
\newline
\verb+http://home.imf.au.dk/matjph/+

\bibitem[H3]{H3} -----, ``Toric varieties,
Hirzebruch surfaces and error-correcting
codes,'' preprint, 2002. 

\bibitem[Han1]{Han1} S. Hansen, ``The geometry of Deligne-Lusztig
varieties: higher-dimensional AG codes,'' PhD thesis, 
Univ. Aarhus, 1999 (advisor J. Hansen).

\bibitem[Han2]{Han2} ------, ``Error-correcting codes from
higher-dimensional varieties,'' Finite Fields and their 
Applications \underline{7}(2001)530-552.

\bibitem[Han3]{Han3} ------, ``Algebraic geometry over finite fields - 
with a view towards
applications,'' Master's thesis, Univ. Aarhus, 1996 (advisor J. Hansen).

\bibitem[J]{J} D. Joyner, MAGMA code: \verb+toric.mag+ code,
\newline
\verb+http://web.usna.navy.mil/~wdj/papers/toric.mag+
\newline
GAP code: \verb+toric.g+ code,
\newline
\verb+http://web.usna.navy.mil/~wdj/papers/toric.g+

\bibitem[JV]{JV} ------ and H. Verrill, ``Notes on toric varieties,''
preprint, available at
\newline
\verb+http://web.usna.navy.mil/~wdj/papers.html+

\bibitem[L1]{L1} J. Little, ``A key equation and the computation of error 
values for codes from order domains'',
\newline
\verb+http://www.arXiv.org/abs/math.AC/0303299+

\bibitem[L2]{L2} ------, ``The Ubiquity of Order Domains for 
the Construction of Error Control Codes,'' 
\newline
\verb+http://www.arXiv.org/abs/math.AC/0304292+

\bibitem[MS]{MS} 
F. MacWilliams and N. Sloane, {\bf The theory of error-correcting codes},
North-Holland, 1977.

\bibitem[MAGMA]{MAGMA} 
W. Bosma, J. Cannon, C. Playoust,
``The MAGMA algebra system, I: The user language,''
J. Symb. Comp., \underline{24}(1997)235-265.
\newline
(See also the MAGMA homepage at
\newline
\verb+ http://www.maths.usyd.edu.au:8000/u/magma/+)

\bibitem[S]{S} H. Stichtenoth, {\bf Algebraic function
fields and codes}, Springer-Verlag, 1993.

\bibitem{VT}[VT] M. Tsfasman,
S. Vladut, {\bf Algebraic-geometric codes}, North-Holland, 1998.

\end{thebibliography}

\end{document}

